% スマホ版専用の組版。本文の文字サイズ・数式のスタイルは本編と同じ。
\setlength{\mathindent}{1zw}
\setlength{\fullwidth}{\textwidth}
\setlength{\columnseprule}{0pt}
\raggedbottom

\makeatletter
% 目次の深い階層でも長い数式を紙幅内に収める。
\renewcommand*{\l@subsection}{\@dottedtocline{2}{2zw}{3.5zw}}
\renewcommand*{\l@subsubsection}{\@dottedtocline{3}{3zw}{4.5zw}}
% 柱は本編と同じ文字・下線・太字のページ番号で、長いものだけ折り返す。
\pagestyle{headings}
\def\@evenhead{\underline{\hbox to\textwidth{\autoxspacing\textbf{\thepage}\hfil
  \parbox[b]{\dimexpr\textwidth-2zw\relax}{\raggedleft\baselineskip=10pt\leftmark}}}}
\def\@oddhead{\underline{\hbox to\textwidth{\autoxspacing
  \parbox[b]{\dimexpr\textwidth-2zw\relax}{\raggedright\baselineskip=10pt\rightmark}\hfil\textbf{\thepage}}}}
\makeatother

% 図の描画内容を変えず、紙幅を超える場合だけ縮小する。
\newsavebox{\BookMobileGraphBox}
\renewcommand{\BookGraph}[2]{%
  \par\smallskip\noindent
  \begin{minipage}{\linewidth}\centering
    \sbox{\BookMobileGraphBox}{%
      \begin{tikzpicture}[scale=0.8,>=Stealth,line cap=round,line join=round,
        every node/.style={font=\small}]#2\end{tikzpicture}}%
    \ifdim\wd\BookMobileGraphBox>\linewidth
      \resizebox{\linewidth}{!}{\usebox{\BookMobileGraphBox}}%
    \else\usebox{\BookMobileGraphBox}\fi
    \BookFigureCaption{#1}%
  \end{minipage}\par\smallskip
}

% 組版ログに入力ファイルを記録して、はみ出しの場所を追えるようにする。
\AddToHook{file/before}{\typeout{MOBILE-BEGIN:\CurrentFile}}
\AddToHook{file/after}{\typeout{MOBILE-END:\CurrentFile}}
\typeout{MOBILE-MATHINDENT:\the\mathindent}
